\documentclass[../../main.tex]{subfiles}
\begin{document}

\section{光滑映射的微分}\label{section:ch03-光滑映射的微分}

为把我们在流形上定义的抽象切空间与 $\mathbb{R}^{n}$ 中的几何切空间联系起来，我们必须考察光滑映射作用于切向量的方式。在欧几里得空间之间的光滑映射情形，该映射在一点处的全导数（由其 Jacobian 矩阵表示）是一个线性映射，代表该映射在给定点附近的"最佳线性逼近"。在流形情形有一个类似的线性映射，但谈论流形之间的线性映射没有意义。取而代之，它将是切空间之间的一个线性映射。

若 $M$ 和 $N$ 是带边或无边的光滑流形，$F:M\to N$ 是一个光滑映射，则对每个 $p\in M$，我们定义一个映射
\begin{align*}
\mathrm{d}F_{p}:T_{p}M\to T_{F(p)}N,
\end{align*}
称为 $F$ 在 $p$ 处的微分（图~\cref{figure:ch03-3.3}），如下所述。给定 $v\in T_{p}M$，我们令 $\mathrm{d}F_{p}(v)$ 为 $F(p)$ 处的导子，它对 $f\in C^{\infty}(N)$ 的作用由规则
\begin{align*}
\mathrm{d}F_{p}(v)(f)=v(f\circ F)
\end{align*}
给出。注意若 $f\in C^{\infty}(N)$，则 $f\circ F\in C^{\infty}(M)$，所以 $v(f\circ F)$ 有意义。算子 $\mathrm{d}F_{p}(v):C^{\infty}(N)\to\mathbb{R}$ 线性，因为 $v$ 线性；它是 $F(p)$ 处的导子，因为对任意 $f,g\in C^{\infty}(N)$ 有
\begin{align*}
\mathrm{d}F_{p}(v)(fg)=v\bigl((fg)\circ F\bigr)=v\bigl((f\circ F)(g\circ F)\bigr)
=f\circ F(p)\,v(g\circ F)+g\circ F(p)\,v(f\circ F)
=f\bigl(F(p)\bigr)\mathrm{d}F_{p}(v)(g)+g\bigl(F(p)\bigr)\mathrm{d}F_{p}(v)(f).
\end{align*}

\begin{figure}[ht]
\centering
\includegraphics[width=0.8\textwidth]{fig-03-3.png}
\caption{微分}\label{figure:ch03-3.3}
\end{figure}

\begin{proposition}[微分的性质]\label{proposition:ch03-微分性质}
设 $M,N,P$ 是带边或无边的光滑流形，$F:M\to N$ 与 $G:N\to P$ 是光滑映射，$p\in M$。
\begin{enumerate}[(a)]
\item $\mathrm{d}F_{p}:T_{p}M\to T_{F(p)}N$ 线性。
\item $\mathrm{d}(G\circ F)_{p}=\mathrm{d}G_{F(p)}\circ\mathrm{d}F_{p}:T_{p}M\to T_{G\circ F(p)}P$。
\item $\mathrm{d}(\mathrm{Id}_{M})_{p}=\mathrm{Id}_{T_{p}M}:T_{p}M\to T_{p}M$。
\item 若 $F$ 是微分同胚，则 $\mathrm{d}F_{p}:T_{p}M\to T_{F(p)}N$ 是同构，且 $(\mathrm{d}F_{p})^{-1}=\mathrm{d}(F^{-1})_{F(p)}$。
\end{enumerate}
\end{proposition}

\begin{exercise}\label{exercise:ch03-3.7}
证明命题~\cref{proposition:ch03-微分性质}。
\end{exercise}

我们对微分的第一个重要应用，将是利用坐标图把流形上一点处的切空间与欧几里得切空间联系起来。但我们必须先解决一个重要的技术问题：虽然切空间是用整个流形上的光滑函数定义的，坐标图一般却只定义在开子集上。关键一点，表现在下一个命题中，是切向量局部地作用。

\begin{proposition}\label{proposition:ch03-切向量局部作用}
设 $M$ 是带边或无边的光滑流形，$p\in M$，$v\in T_{p}M$。若 $f,g\in C^{\infty}(M)$ 在 $p$ 的某个邻域上一致，则 $vf=vg$。
\end{proposition}

\begin{proof}
令 $h=f-g$，于是 $h$ 是一个在 $p$ 的某个邻域上恒为零的光滑函数。令 $\psi\in C^{\infty}(M)$ 为一个光滑鼓包函数，它在 $h$ 的支集上恒等于 $1$，且支集包含在 $M\setminus\{p\}$ 中。因为在 $h$ 非零处 $\psi\equiv 1$，乘积 $\psi h$ 恒等于 $h$。由于 $h(p)=\psi(p)=0$，引理~\cref{lemma:ch03-流形切向量性质} 蕴涵 $vh=v(\psi h)=0$。由线性性，这蕴涵 $vf=vg$。
\end{proof}

利用这个命题，我们可以把一个开子流形的切空间等同于整个流形的切空间。

\begin{proposition}[开子流形的切空间]\label{proposition:ch03-开子流形切空间}
设 $M$ 是带边或无边的光滑流形，$U\subset M$ 是一个开子集，$\iota:U\hookrightarrow M$ 是包含映射。对每个 $p\in U$，微分 $\mathrm{d}\iota_{p}:T_{p}U\to T_{p}M$ 是一个同构。
\end{proposition}

\begin{proof}
为证明单射，设 $v\in T_{p}U$ 且 $\mathrm{d}\iota_{p}(v)=0\in T_{p}M$。令 $B$ 为 $p$ 的一个邻域，使 $\overline{B}\subset U$。若 $f\in C^{\infty}(U)$ 任意，光滑函数延拓引理保证存在 $\widetilde{f}\in C^{\infty}(M)$，使 $\widetilde{f}\equiv f$ 在 $\overline{B}$ 上。于是由于 $f$ 与 $\widetilde{f}|_{U}$ 是 $U$ 上在 $p$ 的一个邻域上一致的光滑函数，命题~\cref{proposition:ch03-切向量局部作用} 蕴涵
\begin{align*}
vf=v\bigl(\widetilde{f}|_{U}\bigr)=v(\widetilde{f}\circ\iota)=\mathrm{d}\iota(v)_{p}\widetilde{f}=0.
\end{align*}
由于这对每个 $f\in C^{\infty}(U)$ 成立，可得 $v=0$，故 $\mathrm{d}\iota_{p}$ 单射。

另一方面，为证明满射，任取 $w\in T_{p}M$。定义一个算子 $v:C^{\infty}(U)\to\mathbb{R}$，令 $vf=w\widetilde{f}$，其中 $\widetilde{f}$ 是 $M$ 上任意在 $\overline{B}$ 上与 $f$ 一致的光滑函数。由命题~\cref{proposition:ch03-切向量局部作用}，$vf$ 与 $\widetilde{f}$ 的选取无关，故 $v$ 良定义，且容易验证它是 $C^{\infty}(U)$ 在 $p$ 处的一个导子。对任意 $g\in C^{\infty}(M)$，
\begin{align*}
\mathrm{d}\iota_{p}(v)g=v(g\circ\iota)=w(\widetilde{g}\circ\iota)=wg,
\end{align*}
其中最后两个等号成立，因为 $g\circ\iota$、$\widetilde{g}\circ\iota$ 和 $g$ 都在 $B$ 上一致。因此，$\mathrm{d}\iota_{p}$ 也满射。
\end{proof}

给定一个开子集 $U\subset M$，$T_{p}U$ 与 $T_{p}M$ 之间的同构 $\mathrm{d}\iota_{p}$ 是典范定义的，独立于任何选择。从今以后，对任意 $p\in U$，我们把 $T_{p}U$ 与 $T_{p}M$ 等同起来。这个等同只不过是如下观察：$\mathrm{d}\iota_{p}(v)$ 与 $v$ 是同一个导子，只是看作作用在更大流形 $M$ 上的函数、而非 $U$ 上的函数上。由于一个导子对一个函数的作用只依赖于该函数在任意小邻域内的值，这是一个无害的等同。特别地，这意味着任意切向量 $v\in T_{p}M$ 可以无歧义地作用于只定义在 $p$ 邻域上的函数，而不必定义在整个 $M$ 上。

\begin{proposition}[切空间的维数]\label{proposition:ch03-切空间维数}
若 $M$ 是一个 $n$ 维光滑流形，则对每个 $p\in M$，切空间 $T_{p}M$ 是一个 $n$ 维向量空间。
\end{proposition}

\begin{proof}
给定 $p\in M$，令 $(U,\varphi)$ 为包含 $p$ 的一个光滑坐标图。由于 $\varphi$ 是从 $U$ 到开子集 $\widetilde{U}\subset\mathbb{R}^{n}$ 的微分同胚，由命题~\cref{proposition:ch03-微分性质}(d)，$\mathrm{d}\varphi_{p}$ 是从 $T_{p}U$ 到 $T_{\varphi(p)}\widetilde{U}$ 的同构。由于命题~\cref{proposition:ch03-开子流形切空间} 保证 $T_{p}M\cong T_{p}U$ 且 $T_{\varphi(p)}\widetilde{U}\cong T_{\varphi(p)}\mathbb{R}^{n}$，可得 $\dim T_{p}M=\dim T_{\varphi(p)}\mathbb{R}^{n}=n$。
\end{proof}

接着我们需要对带边流形证明一个类似结果。事实上，若 $M$ 是一个带边 $n$ 流形，一开始并不清楚人们应当期望 $M$ 边界点处的切空间是什么样子。它应当是一个 $n$ 维向量空间，像内点处的切空间那样？还是应当是 $(n-1)$ 维，像边界那样？又或者它应当是一个 $n$ 维半空间，像 $M$ 局部所建模的空间 $\mathbb{H}^{n}$ 那样？

正如我们下面将证明的，我们的定义蕴涵边界点处的切空间是一个 $n$ 维向量空间（图~\cref{figure:ch03-3.4}），正像内点处的切空间一样。这也许看起来、也许并不像是几何上最直观的选择，但它的优点是使带边流形上大多数几何对象的定义看起来与流形上的完全相同。

首先，我们需要对 $a\in\partial\mathbb{H}^{n}$ 的点，把切空间 $T_{a}\mathbb{H}^{n}$ 与 $T_{a}\mathbb{R}^{n}$ 联系起来。由于 $\mathbb{H}^{n}$ 不是 $\mathbb{R}^{n}$ 的开子集，命题~\cref{proposition:ch03-开子流形切空间} 不适用。作为替代，我们有下面的引理。

\begin{lemma}\label{lemma:ch03-包含映射微分同构}
令 $\iota:\mathbb{H}^{n}\hookrightarrow\mathbb{R}^{n}$ 表示包含映射。对任意 $a\in\partial\mathbb{H}^{n}$，微分 $\mathrm{d}\iota_{a}:T_{a}\mathbb{H}^{n}\to T_{a}\mathbb{R}^{n}$ 是一个同构。
\end{lemma}

\begin{proof}
设 $a\in\partial\mathbb{H}^{n}$。为证明 $\mathrm{d}\iota_{a}$ 单射，假设 $\mathrm{d}\iota_{a}(v)=0$。设 $f:\mathbb{H}^{n}\to\mathbb{R}$ 光滑，令 $\widetilde{f}$ 为 $f$ 到定义在整个 $\mathbb{R}^{n}$ 上的光滑函数的任意延拓。（由光滑函数延拓引理（引理~2.26），这样的延拓存在。）于是 $\widetilde{f}\circ\iota=f$，故
\begin{align*}
vf=v(\widetilde{f}\circ\iota)=\mathrm{d}\iota_{a}(v)\widetilde{f}=0,
\end{align*}
这蕴涵 $\mathrm{d}\iota_{a}$ 单射。

为证明满射，任取 $w\in T_{a}\mathbb{R}^{n}$。定义 $v\in T_{a}\mathbb{H}^{n}$ 为
\begin{align*}
vf=w\widetilde{f},
\end{align*}
其中 $\widetilde{f}$ 是 $f$ 的任意光滑延拓。用 $T_{a}\mathbb{R}^{n}$ 的标准基把 $w$ 写成 $w=w^{i}\partial/\partial x^{i}|_{a}$，这意味着
\begin{align*}
vf=w^{i}\frac{\partial\widetilde{f}}{\partial x^{i}}(a).
\end{align*}
这与 $\widetilde{f}$ 的选取无关，因为由连续性，$\widetilde{f}$ 在 $a$ 处的各阶导数由 $f$ 在 $\mathbb{H}^{n}$ 中的导数决定。容易验证 $v$ 是 $a$ 处的导子且 $w=\mathrm{d}\iota_{a}(v)$，故 $\mathrm{d}\iota_{a}$ 满射。
\end{proof}

正像当 $U$ 是 $M$ 的开子集时我们用命题~\cref{proposition:ch03-开子流形切空间} 把 $T_{p}U$ 与 $T_{p}M$ 等同起来那样，当 $a\in\partial\mathbb{H}^{n}$ 时我们用这个引理把 $T_{a}\mathbb{H}^{n}$ 与 $T_{a}\mathbb{R}^{n}$ 等同起来，并且我们不在记号上区分 $T_{a}\mathbb{H}^{n}$ 的一个元素与其在 $T_{a}\mathbb{R}^{n}$ 中的像。

\begin{figure}[ht]
\centering
\includegraphics[width=0.8\textwidth]{fig-03-4.png}
\caption{带边流形的切空间}\label{figure:ch03-3.4}
\end{figure}

\begin{proposition}[带边流形上切空间的维数]\label{proposition:ch03-带边切空间维数}
设 $M$ 是一个 $n$ 维光滑带边流形。对每个 $p\in M$，$T_{p}M$ 是一个 $n$ 维向量空间。
\end{proposition}

\begin{proof}
任取 $p\in M$。若 $p$ 是内点，则由于 $\mathrm{Int}\,M$ 是 $M$ 的一个开子流形，命题~\cref{proposition:ch03-开子流形切空间} 蕴涵 $T_{p}(\mathrm{Int}\,M)\cong T_{p}M$。由于 $\mathrm{Int}\,M$ 是一个光滑无边 $n$ 流形，它所有切空间的维数都是 $n$。

另一方面，若 $p\in\partial M$，令 $(U,\varphi)$ 为包含 $p$ 的一个光滑边界坐标图，令 $\widetilde{U}=\varphi(U)\subset\mathbb{H}^{n}$。有同构 $T_{p}M\cong T_{p}U$（由命题~\cref{proposition:ch03-开子流形切空间}）；$T_{p}U\cong T_{\varphi(p)}\widetilde{U}$（由命题~\cref{proposition:ch03-微分性质}(d)，因 $\varphi$ 是微分同胚）；$T_{\varphi(p)}\widetilde{U}\cong T_{\varphi(p)}\mathbb{H}^{n}$（再由命题~\cref{proposition:ch03-开子流形切空间}）；以及 $T_{\varphi(p)}\mathbb{H}^{n}\cong T_{\varphi(p)}\mathbb{R}^{n}$（由引理~\cref{lemma:ch03-包含映射微分同构}）。结论随之得出。
\end{proof}

回顾例~1.24，每个有限维向量空间都有一个自然的光滑流形结构，它独立于任何基或范数的选取。下面的命题表明向量空间的切空间可以自然地等同于该向量空间本身。

设 $V$ 是一个有限维向量空间，$a\in V$。正像我们早先在 $\mathbb{R}^{n}$ 情形做的那样，对任意向量 $v\in V$，我们定义一个映射 $D_{v|_{a}}:C^{\infty}(V)\to\mathbb{R}$ 为
\begin{align}\label{align:ch03-向量空间方向导数}
D_{v|_{a}}f=\left.\frac{\mathrm{d}}{\mathrm{d}t}\right|_{t=0}f(a+tv).
\end{align}

\begin{proposition}[向量空间的切空间]\label{proposition:ch03-向量空间切空间}
设 $V$ 是一个带其标准光滑流形结构的有限维向量空间。对每个点 $a\in V$，由~\eqref{align:ch03-向量空间方向导数} 定义的映射 $v\mapsto D_{v|_{a}}$ 是从 $V$ 到 $T_{a}V$ 的一个典范同构，使得对任意线性映射 $L:V\to W$，下图交换：
\begin{center}
$V\xrightarrow{\cong}T_{a}V$，$W\xrightarrow{\cong}T_{La}W$，前者经 $L$ 与后者经 $\mathrm{d}L_{a}$ 对应。
\end{center}
\end{proposition}

\begin{proof}
一旦我们为 $V$ 选取一组基，就能用与命题~\cref{proposition:ch03-导子与几何切向量} 证明中相同的论证，证明 $D_{v|_{a}}$ 确实是 $a$ 处的导子，且映射 $v\mapsto D_{v|_{a}}$ 是同构。

现在设 $L:V\to W$ 是一个线性映射。由于它关于 $V$ 和 $W$ 任意基选择的分量都是坐标的线性函数，$L$ 光滑。展开定义并利用 $L$ 的线性性，我们计算
\begin{align*}
\mathrm{d}L_{a}\bigl(D_{v|_{a}}\bigr)f=D_{v|_{a}}(f\circ L)
=\left.\frac{\mathrm{d}}{\mathrm{d}t}\right|_{t=0}f\bigl(L(a+tv)\bigr)
=\left.\frac{\mathrm{d}}{\mathrm{d}t}\right|_{t=0}f(La+tLv)
=D_{Lv|_{La}}f.
\end{align*}
\end{proof}

重要的是要理解，每个同构 $V\cong T_{a}V$ 都是典范定义的，独立于任何基的选取（尽管我们使用了一组基的选取来证明它是同构）。由于这个结果，我们可以惯常地把一个有限维向量空间的切向量等同于该空间本身的元素。更一般地，若 $M$ 是一个向量空间 $V$ 的开子流形，我们可以把等同 $T_{p}M\leftrightarrow T_{p}V\leftrightarrow V$ 结合起来，得到 $M$ 的每个切空间与 $V$ 的一个典范等同。例如，由于 $\mathrm{GL}(n,\mathbb{R})$ 是向量空间 $\mathsf{M}(n,\mathbb{R})$ 的一个开子流形，我们可以把它在每一点 $X\in\mathrm{GL}(n,\mathbb{R})$ 处的切空间等同于整个矩阵空间 $\mathsf{M}(n,\mathbb{R})$。

对乘积流形的切空间，还有另一个自然等同。

\begin{proposition}[乘积流形的切空间]\label{proposition:ch03-乘积流形切空间}
设 $M_{1},\dots,M_{k}$ 是光滑流形，对每个 $j$，令 $\pi_{j}:M_{1}\times\cdots\times M_{k}\to M_{j}$ 为到 $M_{j}$ 分量的投影。对任意点 $p=(p_{1},\dots,p_{k})\in M_{1}\times\cdots\times M_{k}$，由
\begin{align}\label{align:ch03-乘积切空间同构}
\alpha(v)=\bigl(\mathrm{d}(\pi_{1})_{p}(v),\dots,\mathrm{d}(\pi_{k})_{p}(v)\bigr)
\end{align}
定义的映射
\begin{align*}
\alpha:T_{p}(M_{1}\times\cdots\times M_{k})\to T_{p_{1}}M_{1}\oplus\cdots\oplus T_{p_{k}}M_{k}
\end{align*}
是一个同构。若其中一个空间 $M_{i}$ 是光滑带边流形，同样成立。
\end{proposition}

\begin{proof}
见习题~3-2。
\end{proof}

再次，由于同构~\eqref{align:ch03-乘积切空间同构} 是典范定义的，独立于任何坐标选择，我们可以把它看作一个典范等同，并且我们将总是这样做。于是，例如，我们把 $T_{(p,q)}(M\times N)$ 与 $T_{p}M\oplus T_{q}N$ 等同，并把 $T_{p}M$ 和 $T_{q}N$ 看作 $T_{(p,q)}(M\times N)$ 的子空间。

\end{document}